\section{Introduction}\label{sec:introduction}
A causal experiment supplies legal preparations, reports, interventions and executable future tests, not a pre-existing parameter manifold. Histories are predictively equivalent when those tests have identical continuation laws and interfaces. After this quotient has been realized, finite observation remains a causal problem: the law available for compression is generated by the observer's own previous acquisitions, while its retained information determines the next acquisition. Neither the dimension of a formal state space nor a terminal codebook settles whether the required information can be retained and used by one finite programme.

There is a further constraint when the physical model is inaccessible. A programme synthesized separately at each possible model is not a programme that can operate without knowing that model. The same action row, report row, stopping rule and readout must work at every occupied parameter and phase. A conditional predictive state may vary with both of these indices as an analyst's certificate, but neither index is automatically a machine register. In particular a parameterwise Bayes action is not a common minimax readout. The mathematical objective is an exact resource--risk criterion with these quantifiers in their operational order, followed by a finite construction that approximates the criterion at the original memory budget.

Theorem~\ref{thm:uniform} is the central result. For a jointly dominated family of prepared compact predictive experiments, it eliminates the common report-to-label kernel through the simultaneous affine inequality \eqref{eq:uniform_dual}. The maximum over next labels comes after both the parameter and phase sums. This gives an exact acquired-law invariant
\[
 R^\Theta_{\aut}(n,M)=\Phi^\Theta(n,M)
     =\inf_{\text{one feasible common flow}}\sup_{\theta\in\Theta}R_\theta.
\]
The invariant retains actual conditional occupancies and predictive barycenters; it is not a formula in ambient dimension. The observer does not receive the true parameter, its posterior profile, or an uncharged clock. The single-law barycentric theory is proved in the later sections as a constituent of this statement, not used to authorize a parameter-dependent programme.

On a common finite-dimensional positive physical realization, Theorem~\ref{thm:uniform_compactness} strengthens this exact criterion: the common-programme optimum is attained, and a strict target below it is excluded by finitely many parameter models. This global statement follows from compactness of the shared report rows through their integrated instrument moments. Reuse of a row creates finite products of these moments, not a discontinuous multiplication of arbitrary pointwise weak-star limits. The finite-model obstruction concerns every programme at the budget, rather than one proposed flow. Without further effective data it supplies neither the number of models nor an algorithm to find them.

The effective clause closes that distinction under explicit uniform presentation hypotheses. A common report reconstruction, support-resolving partition and parameter modulus give finite numbers $L,U$ satisfying
\[
 L-E_h-\beta_Q\le \Phi^\Theta(n,M)\le U+\Omega_n(s),
 \qquad U-L\le\eta.
\]
An explicitly enumerated dyadic programme attains the upper bound at the same $M$ persistent labels. The report error $E_h$, programme/readout resolution $\beta_Q$, parameter-net error $\Omega_n(s)$ and integration tolerance $\eta$ are distinct and computably vanish. This bounds the optimum of \emph{all} Borel programmes at that budget, not merely the best element of a promising finite list. A constant unknown model is used throughout each evaluation. Neither a fresh adversarial model at each call nor the exchange $\inf\sup=\sup\inf$ is invoked. The number of programmes and the offline parameter-net evaluations are explicit; the result is not an efficient planning theorem.

The proof operates on positive unnormalized instruments. Their joint physical-state/label laws have a direct-sum base norm contracted by every common stochastic programme. Averaging a report row against a parameter- and phase-independent reconstruction therefore preserves a common row and gives a uniform error without differentiating a normalized posterior. Support certificates retain the possible legal and stopping paths; dyadic rounding can remove an edge but cannot create a forbidden one. The same norm comparison is uniform over all discontinuous Borel report rows and supplies the parameter modulus needed for global lower as well as upper certificates. This separates predictive-update instability from finite-horizon programme approximation.

Theorems~\ref{thm:uniform_gaussian} and~\ref{thm:uniform_quantum} verify the same theorem directly for two different families of raw experiments. Gaussian observations are untruncated and their transitions may be reducible or deterministic. The quantum family permits noncommuting, unknown channels and pure preparations, while the observer remains a classical finite programme rather than a quantum memory. Corollary~\ref{cor:uniform_singular} treats singular reference laws, controlled atomic tails and flagged changes of rank with their actual masses. The inherited geometry theorems additionally give matching regular and singular resource curves when their sharper mass, covering, curvature or blind-transport assumptions hold. None of those realization conclusions is a premise of the general transfer theorem.

Proposition~\ref{prop:uniform_threshold} supplies a sharp finite test of the new quantifier. A noisy calibration bit is followed by a fair report whose interpretation depends on the fixed unknown bit. The optimal common programme has an exact five-state threshold, whereas a programme furnished with that bit can have zero risk with four states. The lower bound covers randomized updates and the calibration-offset term is attained in the same task. This example distinguishes information patterns; it is not offered as a difficult filtering problem solved by elementary algebra. The earlier stationary-reuse parallelogram and noncontracting acquired-law bottleneck remain separate sharpness witnesses.

Comparison of experiments, convex order, conditional martingale kernels and finite-state-controller synthesis are established subjects\cite{blackwell,strassen,leskela,junges}. Robust synthesis for uncertain POMDPs also predates this work\cite{robustfsc}; hidden-model controller extraction and verification provide a recent computational comparison\cite{hiddenmodel}. The claim here is not priority for robust finite memory, conditional barycenters, weak-star compactness or finite enumeration. It is the exact shared-parameter/shared-phase elimination, the global finite-model consequence in the positive-instrument category, and its uniformly certified same-budget implementation with continuous reports and a fixed physical parameter. Section~\ref{sec:literature} separates these assertions from neighboring finite-controller optimization, model learning, discounted control and zero-delay coding results.

The general passage remains
\[
\begin{gathered}
 \text{prepared experiments}\longrightarrow\text{executable predictive quotients}\longrightarrow\text{causal allocations}\\
 \longrightarrow\text{acquired geometry}\longrightarrow\text{finite-resource risk}.
\end{gathered}
\]
The exact compact category and the effective positive-instrument category are stated separately. The latter requires a known effective family, not knowledge of its true member. No universal data-driven identification rate, arbitrary noncompact quotient theorem, infinite-horizon limit, or jointly optimal programme/workspace/time region is asserted. Programme bits, numerical approximation, physical acquisitions and parameter-independent simulator state are separately charged. These boundaries leave the mother problem intact while specifying the part for which complete transfer and implementation proofs are given.
